\section{实验内容}

\subsection{建立可执行的本地质量门禁}
本题要求复用第 10 题的 greetlab 包，建立一个轻量本地质量检查，详细如下：
\begin{itemize}
    \item 复制 q10 为 q13；在 pyproject.toml 中加入 ruff 和 pytest 的最小配置。
    \item 补充一个正常姓名测试和一个空白姓名测试，每个测试包含有效断言。
    \item 运行 ruff format、ruff check 和 pytest，修复全部问题，不得全局忽略规则。
    \item 编写 check.sh，依次执行格式检查、静态检查和 pytest；运行并保证返回 0。
\end{itemize}

\subsubsection*{配置 ruff 与 pytest}
在 pyproject.toml 中加入 ruff 的规则配置与 pytest 的测试路径配置：
\begin{minted}[frame=single, fontsize=\small]{ini}
[build-system]
requires = ["setuptools>=68"]
build-backend = "setuptools.build_meta"
[project]
name = "greetlab-25020007136"
version = "0.1.0"
[project.scripts]
sdt-greet = "greetlab.cli:main"
[tool.ruff]
line-length = 88
target-version = "py313"
[tool.ruff.lint]
select = ["E", "F", "I", "W"]
[tool.pytest.ini_options]
testpaths = ["."]
python_files = "test_*.py"
\end{minted}

\begin{figure}[H]
    \centering
    \includegraphics[width=0.95\linewidth]{q13_0}
\end{figure}

\subsubsection*{补充测试}
补充正常姓名与空白姓名两个测试，每个测试都包含有效断言：
\begin{minted}[frame=single, fontsize=\small]{python}
import os
import sys

# 将 q13/src 目录加入到 Python 搜索路径中
sys.path.insert(0, os.path.abspath(os.path.join(os.path.dirname(__file__), "src")))

import pytest

from greetlab.cli import main


def test_blank_name_exits_with_2(monkeypatch):
    # 模拟命令行输入
    monkeypatch.setattr(sys, "argv", ["sdt-greet", "--name", "   "])
    with pytest.raises(SystemExit) as exc_info:
        main()
    # 验证退出码为 2
    assert exc_info.value.code == 2


def test_normal_name(monkeypatch, capsys):
    """测试正常姓名，应输出 Hello 问候语并正常退出。"""
    monkeypatch.setattr(sys, "argv", ["sdt-greet", "--name", "Alice"])
    main()

    captured = capsys.readouterr()
    assert "Hello, Alice!" in captured.out


def test_blank_name(monkeypatch, capsys):
    """测试纯空白姓名，应以 SystemExit(2) 结束。"""
    monkeypatch.setattr(sys, "argv", ["sdt-greet", "--name", "   "])
    with pytest.raises(SystemExit) as exc_info:
        main()

    assert exc_info.value.code == 2

    captured = capsys.readouterr()
    assert "error" in captured.err.lower()
\end{minted}

\subsubsection*{编写 check.sh 并运行}
编写 check.sh 依次执行格式检查、静态检查与测试，并保证返回 0：
\begin{minted}[frame=single, fontsize=\small]{bash}
set -e

ruff format
ruff check --fix
pytest
\end{minted}

运行结果如下：
\begin{figure}[H]
    \centering
    \includegraphics[width=0.95\linewidth]{q13_1}
\end{figure}

\begin{figure}[H]
    \centering
    \includegraphics[width=0.95\linewidth]{q13_2}
\end{figure}

\subsection{让 Make 只重建真正受影响的产物}
本题要求编写 Makefile，只重建真正受影响的产物，详细如下：
\begin{itemize}
    \item 创建给定的 data.csv、stats.py、build\_report.py 和 report.md。
    \item 编写 Makefile，包含 all、stats.txt、report.txt 和 clean；准确列出数据文件、脚本和中间产物依赖，并把 all、clean 声明为 .PHONY。
    \item 验证首次 make 生成两个产物；第二次无改动时不执行配方。
    \item touch data.csv 后再次 make，确认重新生成；clean 只删除生成物。
\end{itemize}

给定内容：
\begin{minted}[frame=single, fontsize=\small]{python}
# data.csv
name,value
alpha,2
beta,3
gamma,5

# stats.py
import csv
with open("data.csv") as f:
    total = sum(int(r["value"]) for r in csv.DictReader(f))
open("stats.txt", "w").write(str(total))

# build_report.py
text = open("report.md").read()
total = open("stats.txt").read()
open("report.txt", "w").write(f"{text}\nTotal: {total}\n")

# report.md
# Data Report
\end{minted}

\subsubsection*{编写 Makefile}
将 all、clean 声明为 .PHONY，并为两个产物准确列出依赖：
\begin{minted}[frame=single, fontsize=\small]{make}
# 声明伪目标，防止目录下有同名文件导致 make 误判
.PHONY: all clean

# 默认目标：依赖于两个最终产物
all: stats.txt report.txt

# stats.txt 依赖于源数据 data.csv 和脚本 stats.py
stats.txt: data.csv stats.py
	python3 stats.py

# report.txt 依赖于中间产物 stats.txt、脚本 build_report.py 以及模板文件 report.md
report.txt: stats.txt build_report.py report.md
	python3 build_report.py

# 清理规则：只删除生成的产物，不删除源文件
clean:
	rm -f stats.txt report.txt
\end{minted}

\begin{figure}[H]
    \centering
    \includegraphics[width=0.95\linewidth]{q14}
\end{figure}

\subsection{把本地 API 数据转换为可读报告}
本题要求创建本地 HTTP 服务，再用 curl 与 jq 生成 Markdown 报告，详细如下：
\begin{itemize}
    \item 把给定 JSON 保存为 packages.json，并在 q15 目录运行 python -m http.server 8000。
    \item 使用 curl -fsS 获取 http://127.0.0.1:8000/packages.json。
    \item 使用 jq 筛选 status 为 active 且 downloads 不少于 100 的记录，并按 downloads 降序、name 升序排列。
    \item 编写 api\_report.sh，把结果生成 summary.md，包含标题和 name/version/downloads 三列表格。
\end{itemize}

给定 JSON：
\begin{minted}[frame=single, fontsize=\small]{json}
[
  {"name":"alpha","status":"active","downloads":120,"version":"1.2.0"},
  {"name":"beta","status":"inactive","downloads":900,"version":"2.0.0"},
  {"name":"gamma","status":"active","downloads":450,"version":"1.5.1"},
  {"name":"delta","status":"active","downloads":450,"version":"0.9.0"},
  {"name":"epsilon","status":"active","downloads":80,"version":"3.1.0"}
]
\end{minted}

\subsubsection*{启动本地服务并获取数据}
启动 HTTP 服务，再用 curl 获取数据：
\begin{minted}[frame=single, fontsize=\small]{bash}
cd q15
python3 -m http.server 8000 &
curl -fsS http://127.0.0.1:8000/packages.json
\end{minted}

\subsubsection*{编写 api\_report.sh}
使用 jq 筛选并排序，生成 summary.md：
\begin{minted}[frame=single, fontsize=\small]{bash}
#!/usr/bin/env bash

URL="http://127.0.0.1:8000/packages.json"
OUTPUT="summary.md"

echo "# API Package Report" > "$OUTPUT"
echo "" >> "$OUTPUT"
echo "| name | version | downloads |" >> "$OUTPUT"
echo "|------|---------|-----------|" >> "$OUTPUT"

curl -fsS "$URL" | jq -r '[.[] | select(.status == "active" and .downloads >= 100)] | sort_by([-.downloads, .name])[] | "| \(.name) | \(.version) | \(.downloads) |"' >> "$OUTPUT"
\end{minted}

生成的 summary.md 内容如下：
\begin{verbatim}
# API Package Report

| name | version | downloads |
|------|---------|-----------|
| delta | 0.9.0 | 450 |
| gamma | 1.5.1 | 450 |
| alpha | 1.2.0 | 120 |
\end{verbatim}

\begin{figure}[H]
    \centering
    \includegraphics[width=0.95\linewidth]{q15}
\end{figure}

\subsection{修复并交付一个陌生的小型工具仓库}
本题要求复用第 13 题项目，完成一次综合检查，详细如下：
\begin{itemize}
    \item 复制 q13 为 q16 并初始化/继续使用 Git。
    \item 把问候语实现临时改成始终返回字面量 Hello, name!，运行 check.sh 确认测试失败；随后定位并修复。
    \item 编写最小 Makefile，包含 check、build 和 clean；check 调用 check.sh，build 生成 wheel。
    \item 运行 make check 和 make build，计算 wheel 的 SHA-256，并把最终改动提交为一个内容聚焦的 Git 提交。
\end{itemize}

\subsubsection*{初始化仓库并制造失败}
复制 q13 为 q16，初始化 Git，把实现临时改成始终返回字面量，运行 check.sh 确认测试失败：
\begin{minted}[frame=single, fontsize=\small]{bash}
cp -r q13 q16 && cd q16
git init
./check.sh   # 临时改坏后应失败
\end{minted}

\begin{figure}[H]
    \centering
    \includegraphics[width=0.95\linewidth]{q16_0}
\end{figure}

定位并恢复 cli.py 中的空白校验逻辑后，make check 通过：
\begin{figure}[H]
\centering
\begin{subfigure}[b]{0.49\textwidth}
    \includegraphics[width=\linewidth]{q16_1}
\end{subfigure}
\hfill
\begin{subfigure}[b]{0.49\textwidth}
    \includegraphics[width=\linewidth]{q16_2}
\end{subfigure}
\end{figure}

\subsubsection*{编写 Makefile 并构建}
编写最小 Makefile，check 调用 check.sh，build 生成 wheel：
\begin{minted}[frame=single, fontsize=\small]{make}
.PHONY: check build clean

check:
	./check.sh

build:
	python3 -m build

clean:
	rm -rf dist/ build/ *.egg-info .pytest_cache
	find . -type d -name "__pycache__" -exec rm -rf {} +
\end{minted}

\begin{minted}[frame=single, fontsize=\small]{bash}
make check
make build
\end{minted}

\begin{figure}[H]
    \centering
    \includegraphics[width=0.95\linewidth]{q16_3}
\end{figure}

\subsubsection*{校验和并提交}
计算 wheel 的 SHA-256，并把最终改动提交为内容聚焦的 Git 提交：
\begin{minted}[frame=single, fontsize=\small]{bash}
sha256sum dist/greetlab_25020007136-0.1.0-py3-none-any.whl
git add -A
git commit -m "fix: 恢复空白姓名校验并补齐质量门禁"
\end{minted}

\begin{figure}[H]
\centering
\begin{subfigure}[b]{0.49\textwidth}
    \includegraphics[width=\linewidth]{q16_4}
\end{subfigure}
\hfill
\begin{subfigure}[b]{0.49\textwidth}
    \includegraphics[width=\linewidth]{q16_5}
\end{subfigure}
\end{figure}